\section{Diskussion}
\label{sec:discussion}

Die Forschungsfrage lautete, ob sich Schreibaktivität allein aus den
IMU-Daten einer Smartwatch erkennen lässt, unabhängig von der Person,
die sie trägt. Die Antwort ist ja, mit einer klar benennbaren Grenze. Auf 32
Personen, von denen jede für ihr eigenes Testergebnis nie im Training war,
erreicht ein zeitliches Faltungsnetz auf 5-s-Fenstern 92\,\% Accuracy, und
ein Random Forest auf handgebauten Merkmalen mit kausalem HMM-Filter 90\,\%.
Beide Zahlen sind ohne Kenntnis der Person, ohne Kalibrierungsphase und ohne
Blick in die Zukunft erreichbar. Was fehlt, sind die restlichen acht bis zehn
Prozent, und die Ergebnisse sagen ziemlich genau, wo sie sitzen.

\paragraph{Die Grenze des Merkmalsmodells liegt vermutlich in der
Repräsentation.} Über den größten Teil des Projekts sah es so aus, als läge
die Leistungsgrenze im Signal. Random Forest, MiniRocket, ein vortrainiertes
Modell und die Netze auf 1-s-Fenstern trafen dieselbe Marke und scheiterten
an denselben Personen. Kein Merkmalseingriff, keine Augmentation und keine
Umgewichtung bewegte sie. Der Schritt von 22 auf 32 Personen trennt diese
Lesart, wenn auch nicht sauber: Der Random Forest blieb stehen, die Netze
stiegen — zugleich änderten sich Kohortenzusammensetzung und Protokollanteil,
sodass die Beobachtung nicht allein der Personenzahl zuzuschreiben ist. Die
naheliegende Deutung ist, dass nicht das Signal erschöpft ist, sondern das,
was Statistiken über ein 1-s-Fenster ausdrücken können. Ein Netz, das die
Reihenfolge der Bewegungen innerhalb von fünf Sekunden sieht, profitiert
anders als der Random Forest weiter von zusätzlichen Personen. Dass die Netze im Modern-Pool auf
23 Personen den Schwerkraftvektor nutzen können, während der Random Forest
auf vier Personen daran scheiterte, passt in dasselbe Bild.

\paragraph{Zeitkontext ist der entscheidende Hebel, und man kann ihn nur
einmal einsammeln.} Drei ganz verschiedene Wege führen auf dieselbe Höhe:
ein Netz auf nativen 5-s-Fenstern, ein 1-s-Modell mit HMM, und die Fusion
von Random Forest und Netz. Das HMM hilft nur Modellen ohne eigenes
Gedächtnis. Sobald das Basismodell fünf Sekunden Kontext trägt, überglättet
es. Dass der Rolling-Mean auf keiner Skala hilft, das HMM aber auf jeder,
zeigt, dass es nicht um Mittelung geht, sondern um ein Modell der
Übergänge: Schreiben ist ein Zustand mit Trägheit, und ein Detektor, der das
weiß, darf bei schwacher Evidenz beim letzten Zustand bleiben. Für den
Einsatz ist damit die Frage nach der Latenz entscheidend. Wer eine Anzeige
im Moment will, bekommt 0,90. Wer abends den Tag rekonstruiert, bekommt
0,92, weil der Smoother beide Seiten eines Übergangs sehen darf. Diese
Differenz von knapp zwei Prozentpunkten ist in 31 von 32 Personen dieselbe
und war bei 22 Personen schon vorhanden. Sie gehört zu den robustesten
Befunden der Arbeit.

\paragraph{Der Restfehler ist personen- und aufgabenspezifisch, nicht
diffus.} Die Fehler des Modells liegen nicht gleichmäßig über allen
Nicht-Schreib-Situationen, sondern konzentrieren sich auf zwei Muster: Tippen
auf Tastatur oder Handy bei einem Teil der Personen und sehr leises oder
stark unterbrochenes Schreiben bei einem anderen Teil. Beide Muster sind am
Handgelenk physikalisch begründet. Wer beim Tippen die Hand ähnlich bewegt
wie beim Schreiben, erzeugt ähnliche Rotationen. Wer sehr leicht schreibt,
erzeugt kaum mehr Bewegung als in Ruhe. Für das Tippen konnten wir
zeigen, dass die trennende Information in den Merkmalen vorhanden ist, aber
von der Mehrheit der Merkmale überstimmt wird, und dass alle Versuche, sie
per Merkmalsdesign oder Gewichtung durchzusetzen, gescheitert sind. Der
gepoolte Tipp-Fehler ist mit den zusätzlichen Personen von 36 auf 22
Prozent gefallen, ohne dass jemand daran gearbeitet hätte. Mehr Personen mit
verschiedenen Tippstilen sind offenbar der Hebel, den Merkmalsdesign nicht
ersetzen kann.

\paragraph{Was die Zahlen messen, entscheidet die Labeldefinition.} Das
Schließen von Lücken bis 2,5~s macht aus der Frage \emph{berührt der Stift
das Papier} die Frage \emph{ist die Person im Schreibmodus}. Für einen
Schreibzeit-Tracker ist das die richtige Frage, und sie ist für das Modell
lernbar, weil das Handgelenk in Mikropausen dasselbe tut wie in den Strichen.
Die Kehrseite ist sichtbar: Rechenaufgaben mit Denkpausen über 2,5~s
bleiben schwierig, weil die Semantik dort ausfranst, und ein Schätzer für den
Schreibanteil überschätzt die rohe Stiftzeit systematisch um rund 20
Prozentpunkte. Beides ist kein Modellfehler, sondern eine Folge der
Definition, und beides muss bei jeder Anwendung mitgedacht werden.

\paragraph{Wo die Ergebnisse tragen und wo sie Hypothesen bleiben.} Belegt
sind, jeweils gepaart über Personen und mit $p < 0{,}05$: der Gewinn des
HMM-Filters, der Vorsprung des Smoothers, der Vorsprung der nativen
5-s-Netze gegenüber dem Random Forest bei 15 und 20 Personen, der Gewinn der
Fusion, die Kosten des fehlenden Gyroskops und die Wirkung der Korrekturen
an der Zeitachse. Nominal, aber nicht belegt sind die Rangfolge der beiden
Netze, der Vorsprung der Netze bei 32 Personen (der Abstand ist groß, die
personenweisen Vorhersagen liegen aber nicht vor) und die Wirkung des
Schwerkraftvektors auf tcn6. Hypothesen sind die Erklärungen: dass der
rekurrente Kopf Lageinformation aus dem Verlauf rekonstruiert, und dass die
Tipp-Verwechslung mit mehr Personen weiter schrumpft. Wir haben im Verlauf
des Projekts zweimal einen scheinbar sauberen Befund zurücknehmen müssen,
einmal wegen eines Sortierfehlers, einmal wegen eines verschobenen
Zeitversatzes. Beide Male war nicht die Statistik falsch, sondern die Daten
darunter. Das ist der Grund, warum diese Arbeit so viel Platz auf
Zeitachsen und Kontrollen verwendet.
